\documentclass[a4paper,11pt]{article}

\usepackage{amsmath,amssymb, amsthm}
\usepackage{charter}
%\usepackage{calc}

\usepackage{fullpage}

\usepackage{multirow}
%\usepackage[cp1250]{inputenc}
%\usepackage[T1]{fontenc}
%\usepackage{calligra}
%\usepackage[slovak]{babel}
\usepackage{amsfonts}
\usepackage{layout}
\usepackage[dvips]{graphicx}
\usepackage{url}
\usepackage{color}
\usepackage{graphicx}
\usepackage{algorithmic}
\usepackage{algorithm}
\usepackage{verbatim}
\usepackage{epsfig}
\usepackage[colorlinks,
            linkcolor=red,
            anchorcolor=blue,
            citecolor=blue
            ]{hyperref}
\usepackage{xcolor}

\renewcommand{\algorithmicrequire}{\textbf{Input:}}
\renewcommand{\algorithmicensure}{\textbf{Output:}}

\newcommand{\hlight}[2]{\noindent\colorbox{#1}{%
    \parbox{\dimexpr\linewidth-1\fboxsep}% a box with line-breaks that's just wide enough
        {#2%
        }}}

\newcommand{\HRule}{\rule{\linewidth}{0.5mm}}
 \newcommand{\xtilde}{\tilde x}
 \newcommand{\vtilde}{\tilde v}
 \newcommand{\Embb}{\E}
 \newcommand{\xbar}{\bar x}
 \newcommand{\ie}{\textit{i.e.}}
 
\newcommand{\main}[1]{\footnotesize\textbf{#1}} 
\newcommand{\todo}[1]{{\color{red}#1}}
\newcommand\tagthis{\addtocounter{equation}{1}\tag{\theequation}}


\newcommand{\del}[1]{{\color{red}#1}}

\let\la=\langle
\let\ra=\rangle

\newcommand{\st}{\;:\;}
\newcommand{\ve}[2]{\langle #1 ,  #2 \rangle}

\newcommand{\eqdef}{\stackrel{\text{def}}{=}}

\newcommand{\ii}{{}^{(i)}}

\newcommand{\R}{\mathbb{R}}
\newcommand{\Prob}{\mathbf{Prob}}
\newcommand{\E}{\mathbb{E}}
\newcommand{\Q}{\mathbb{Q}}
\newcommand{\Z}{\mathbb{Z}}

\newcommand{\vc}[2]{#1^{(#2)}}
\newcommand{\nc}[2]{{\color{red}\|#1\|_{(#2)}}}
\newcommand{\ncs}[2]{\|#1\|^2_{(#2)}}
\newcommand{\ncc}[2]{{\color{red}\|#1\|^*_{(#2)}}}
\newcommand{\ls}[1]{{\color{red} \mathcal S(#1)}}
\newcommand{\Rw}[2]{\mathcal R_{#1}(#2)}
\newcommand{\Rws}[2]{\mathcal R^2_{#1}(#2)}
\newcommand{\m}[1]{~\mbox{#1}~}

\newcommand{\nbp}[2]{\|#1\|_{(#2)}}   % norm block primal
\newcommand{\nbd}[2]{\|#1\|_{(#2)}^*} % norm block dual

\newcommand{\lf}{\mathcal L}
\newcommand{\U}{U}
\newcommand{\N}{N}
\newcommand{\mc}[1]{\mathcal #1}

\newcommand{\mLi}{{\color{red}m^{(i)}}}
\newcommand{\gLi}{{\color{red}g^{(i)}}}
%\newcommand{\TLi}{{\color{red}T_L^{(i)}}}
\newcommand{\TLi}[1]{{\color{blue}T^{(#1)}}}

\newcommand{\Lip}{L}

\newcommand{\Rc}[1]{{\color{red}  \mathbf{RC}_{(#1)}}}
\newcommand{\NRCDM}{{\color{red}NRCDM}\  }
\newcommand{\nnz}[1]{{\color{red}\|#1\|_0}}
% sets
\DeclareMathOperator{\card}{card}       % cardinality of a set
\DeclareMathOperator{\diam}{diam}       % diameter of a set
\DeclareMathOperator{\MVEE}{MVEE}       % minim volume enclosing ellipsoid of a set
\DeclareMathOperator{\vol}{vol}         % volume of a set

\DeclareMathOperator{\prox}{prox}         

% statistical
\DeclareMathOperator{\Exp}{\mathbf{E}}           % expectation
\DeclareMathOperator{\Cov}{Cov}         % covariance
\DeclareMathOperator{\Var}{Var}         % variance
\DeclareMathOperator{\Corr}{Corr}       % correlation

% functions and operators
\DeclareMathOperator{\signum}{sign}     % signum/sign of a scalar
\DeclareMathOperator{\dom}{dom}         % domain
\DeclareMathOperator{\epi}{epi}         % epigraph
\DeclareMathOperator{\Ker}{null}        % nullspace/kernel
\DeclareMathOperator{\nullspace}{null}  % nullpsace
\DeclareMathOperator{\range}{range}     % range
\DeclareMathOperator{\Image}{Im}        % image
\DeclareMathOperator{\argmin}{argmin}        % argmin

% topology
\DeclareMathOperator{\interior}{int}    % interior
\DeclareMathOperator{\ri}{rint}         % relative interior
\DeclareMathOperator{\rint}{rint}       % relative interior
\DeclareMathOperator{\bdry}{bdry}       % boundary
\DeclareMathOperator{\cl}{cl}           % closure

% vectors, matrices
\DeclareMathOperator{\linspan}{span}
\DeclareMathOperator{\linspace}{linspace}
\DeclareMathOperator{\cone}{cone}

\DeclareMathOperator{\tr}{tr}           % trace
\DeclareMathOperator{\rank}{rank}       % rank
\DeclareMathOperator{\conv}{conv}       % convex hull
\DeclareMathOperator{\Diag}{Diag}       % Diag(v) = diagonal matrix with v_i on the diagonal
\DeclareMathOperator{\diag}{diag}       % diag(D) = the diagonal vector of matrix D

\DeclareMathOperator{\Arg}{Arg}         % Argument

%\renewcommand{\qedsymbol}{\ding{114}}


\newtheorem{assumption}{Assumption}
\newtheorem{lemma}{Lemma}
\newtheorem{algorithms}{Algorithm}
\newtheorem{theorem}{Theorem}
\newtheorem{proposition}{Proposition}
\newtheorem{example}{Example}
\newtheorem{remark}{Remark}

\theoremstyle{plain}

\newtheorem{prop}[theorem]{Proposition}
\newtheorem{cor}[theorem]{Corollary}
\newtheorem{lem}[theorem]{Lemma}
\newtheorem{claim}[theorem]{Claim}
%\newtheorem{remark}[theorem]{Remark}

\theoremstyle{definition}

\newtheorem{exercise}[theorem]{Exercise}

\newtheorem{rem}[theorem]{Remark}
\newtheorem{que}[theorem]{Question}
\newtheorem{definition}[theorem]{Definition}

 
\begin{document}

\begin{titlepage}
\begin{center}

% Upper part of the page. The '~' is needed because \\
% only works if a paragraph has started.

\includegraphics[width=0.15\textwidth]{lehigh.jpg}~\\[1cm]

\textsc{\LARGE Lehigh University}\\[1.5cm]

\textsc{\Large 418 Final Project}\\[0.5cm]

% Title
\HRule \\[0.4cm]
{ \Large \bfseries On the exact separation of mixed integer knapsack cuts\\[0.4cm] }
%{ \small \bfseries Based on \cite{csizmadia2012s} and some related papers \\[0.4cm] }

\HRule \\[1.5cm]

% Author and supervisor
\noindent
\begin{minipage}{0.4\textwidth}
\begin{flushleft} \large
\emph{Author:}\\
Chenxin Ma, Xi He
\end{flushleft}
\end{minipage}%
\begin{minipage}{0.4\textwidth}
\begin{flushright} \large
\emph{Supervisor:} \\
Prof.~Ted Ralphs
\end{flushright}
\end{minipage}

\vfill

% Bottom of the page
{\Large Fall 2014}

\end{center}
\end{titlepage}

\tableofcontents
\newpage

\section{Introduction} 
Consider a positive integer $n$, letting $b\in\Q$, $a\in \Q^n$, $l\in \{\Q\cup \{-\infty\}\}^n$, $u\in \{\Q\cup \{+\infty\}\}^n$ and $I\subset N :=
\{1,...,n\}$. The Mixed Integer Knapsack Set is defined as
\begin{equation}\label{K}
 K=\{x\in \R^n : ax\leq b, l\leq x\leq u, x_i\in \Z, \forall i\in I\}.
\end{equation}

Furthermore, if we have $c\in \Q^n$ and assume $l_i$ is finite for each $i\in N$, then the Mixed Integer Knapsack Problem (MIKP) can be 
described as:
\begin{equation}\label{MIKP}
 \max \{cx:x\in K\}
\end{equation}

%\section{Objectives}

In this project, we concern two algorithms developed from \cite{fukasawa2011exact}. The first one is to derive strongly valid inequalities for mixed integer knapsack set \eqref{K}, \ie, for a given 
point $x^\star\in \Q^n$, we focus on
\begin{itemize}
 \item identifying an inequality $(\pi,\pi_0)$ which is valid for $\conv(K)$ (equivalent valid for $K$) and violated by $x^\star$.
 \item proving no such inequality exists
\end{itemize}

This is of great practical importance since by obtaining valid inequalities for \eqref{K}, we can also obtain valid inequalities for general Mixed Integer Programming 
(MIP) problems. Suppose $D\in \Q^{m\times n}$ and $d\in \Q^n$, the mixed integer set can be described as
\begin{equation}\label{P}
 P=\{x\in \R^n: Dx\leq d, l\leq x\leq u, x_i\in \Z, \forall i \in I\}.
\end{equation}

Then if $(a,b)$ in \eqref{K} can be obtained as a nonnegative linear combination of row from $(D,d)$, then $P\in K$ and we can find that any inequality
 which is valid for  $K$ will also be valid for $P$. Actually, the valid inequality is called as $knapsack~inequality$ or $knapsack~cut$ of $P$.
 
The basic idea is to transform the problem of getting a valid inequality to a linear separation problem:
\begin{align}
LP_1: \min &\sum_{i=1}^{n} u_i+v_i \notag\\
s.t. \;  & \pi x^k-\pi_0\leq 0,  \quad \forall k\in\{1,2,..,q\} \notag\\
& \pi x^r\leq 0, \quad~ \qquad \forall  k\in\{1,2,..,t\} \notag\\
& \pi x^*-\pi_0 = 1 \notag\\
& \pi+u-v = 0 \notag\\
&u\geq 0, v\geq 0\notag,
\end{align}
where $\{x^1, x^2,...,x^q\}$ and $\{r^1,r^2,..., r^l\}$ represent the extreme points and rays of $conv(K)$. The following steps can be applied to speed up the procedure
\begin{itemize}
 \item Eliminating variables from $LP_1$
 \item Fixing variables and lifting back again
\end{itemize}

The second one is an efficient branch and bound algorithm with both domination branching and reduced-cost bound improvement (KBB) for the Mixed 
Knapsack Integer Problem \eqref{MIKP}. Most modern algorithms for solving Knapsack Problem are based either on branch-and-bound or on dynamic 
programming. However, the most efficient codes seldom make explicit use of Linear Programming. Here we will develop an LP-based branch and bound 
approach, and three strategies are used to speedup solving the problem:
\begin{itemize}
 \item Detecting unbounded solutions and Preprocessing,
 \item Branch and bound (depth-first-search strategy), by using variable reduced-cost information to improve variable bounds at each node of the tree,
 \item Domination (To use the lexicographic and cost dominance to significantly reducing the search space). 
\end{itemize}

The algorithm exploits dominance conditions, which differs from traditional linear-programming based algorithms by allowing feasible 
solutions to be pruned during the branching phase. The idea is that feasible solutions will only be pruned if either (a) they are not 
optimal (cost-domination-criteria), or (b) if they are optimal, but somewhere else in the tree it is known that there is another optimal 
solution. In the computational study, we will try to see how  domination plays an role in this algorithm. In fact, according to the literature, the 
proposed algorithm performs quite well in practice, outperforming the general-use mixed integer programming solver CPLEX. 

%In this project, we will focus on cutting plane for large Mixed Integer Programs (MIPs) in the context of Branch and Cut algorithms, 
%which is studied in \cite{fukasawa2011exact}. More specifically, we will study an efficient algorithm (KBB) for the Mixed Integer Knapsack Program 
%(MIKP). 

\section{Research plan}
We will implement the algorithm in \cite{fukasawa2011exact}, which contains:
\begin{itemize}
\item a simple eight-step pre-processing procedure in order to reduce an instance of MIKP to another, equivalent
instance of MIKP which is easier to solve.
\item a depth-first-search branch and bound algorithm which always branches on the unique fractional variable. We use a simple 
linear programming algorithm, a variation of Dantzig’s algorithm \cite{dantzig1957discrete}, which runs in linear time by taking advantage 
of the fact that variables are sorted by decreasing efficiency.
\item using dominance to improve the branch and bound search. In fact, lexicographic and cost dominance allow us to disregard feasible 
solutions that are not the unique lexicographically smallest optimum solution, hence significantly reducing the search space.
\end{itemize}
After implementing it, we will conduct extensive computational experiments. The goal is to reproduce the improvement that the new algorithm 
brings in \cite{fukasawa2011exact}. Moreover, we will try to combine different branching strategies (see \cite{linderoth2005noncommercial}) and 
bound improvement to compare the different versions of the algorithm. If possible, we will take a glimpse to use this knapsack cut generation 
procedure to solve some classes of mixed integer programming problems.

\bibliographystyle{plain} 
\bibliography{literature}
 
\end{document}
